\subsection{Deviations from the Wireframes}\label{subsect:Deviations}

The prototype remained the visual reference during implementation. Some changes became necessary after connecting the screens to real data and testing them on mobile devices.

\begin{itemize}
    \item \textbf{One professional role:} the early material described Personal Trainers, nutritionists, and dietitians as separate administrative experiences. The implementation uses one professional role with a specialty. This keeps one interface and enables or hides actions according to the professional's profile;
    \item \textbf{Profile access:} Profile was moved from the bottom navigation to the avatar in the top bar. This leaves four primary destinations per role and avoids an overcrowded navigation bar;
    \item \textbf{Contextual booking:} appointment creation is presented in a bottom sheet rather than as a separate route. The selected professional and day remain visible in the background, reducing the risk of losing context;
    \item \textbf{Data-backed information:} fields that do not exist in the database are not invented for presentation. The professional biography replaces the fixed experience value from the prototype, and the membership identifier is derived from the user's real profile identifier;
    \item \textbf{Access scanner:} the professional QR scanner and its top-bar entry point were added after the first wireframes. They are required to validate the rotating badge designed for the member;
    \item \textbf{Weekly workout runs:} the original live-session wireframes represented a workout that could be completed once. The implemented design records a new run whenever a session is trained, derives progress from the current week, and preserves previous results. It also adds a persistent mini-player, end-of-run confirmation, and a run-specific summary;
    \item \textbf{Draft-then-commit plan creation:} the initial workout form created only one session before returning to the plan. The final wizard keeps any number of sessions in a local draft, adds a review step, and writes the complete plan only after confirmation. The nutrition editor follows the same pattern with day types and meals;
    \item \textbf{Workout-plan history:} editing produces a replacement plan instead of deleting or overwriting prescriptions already referenced by completed workouts. The latest version is shown as active, while previous versions remain available to the database history;
    \item \textbf{Nutrition day types:} the prototype showed the same meals independently of the selected date. The implemented plan assigns reusable day types to weekdays, so the week selector now changes the meals displayed and each meal can include macronutrients and alternatives;
    \item \textbf{No PDF export of the nutrition plan:} the prototype's client screen offered a \textit{View Full PDF Plan} action beside the assigned diet. The implemented plan is structured data the application reads directly, so a generated document would be a second copy of the same information, free to drift from the rows once a day type or a meal is corrected, and it would require file storage the project does not otherwise need. The plan is therefore consulted in the interface, where it is always the current version;
    \item \textbf{Membership enforcement:} subscription status was originally presented as information. In the implemented flow it also controls access to booking, plan creation, gym entry, and reward earning. A professional can suspend or reactivate an assigned client, while the member receives a persistent explanation when services are paused;
    \item \textbf{Body measurements removed:} the progress area was designed around periodic body measurements entered by hand. Nothing in the application produced that data, and asking a professional to type it during a session was slower than the paper card it replaced. Progress now reports what the app actually records: weekly training, the detail of each run, and the member's closing note;
    \item \textbf{Exercise media:} the prototype's graphic exercise blocks were retained as placeholders. Producing and licensing a complete set of images or videos was outside the project scope, and using inconsistent external material would have reduced the visual coherence of the application.
\end{itemize}

These deviations preserve the purpose of the original screens while adapting them to the implemented flows.

A last group of changes came from the prototype being right on the artboard and wrong on the phone, which is worth recording because it is a limit of static prototyping rather than a design mistake. The empty circle drawn on appointment cards was meant as a quick "mark as done" control; that control was never built, and the circle survived as decoration that looks like a checkbox. The calendar was designed in a collapsed and an expanded state, which sit side by side in Figma but force the running application to open in whichever of the two is less useful. The row of dots representing the weekly goal is drawable for a fixed number of sessions and runs off the screen at six or more. None of the three is visible until the screen is real, on a real display, holding real data.
